\documentclass[12pt,oneside]{book}

% =========================================================
% METADATOS
% =========================================================
\newcommand{\universidad}{Universidad de Buenos Aires (UBA)}
\newcommand{\facultad}{Facultad de Derecho}
\newcommand{\materia}{Derecho Procesal Civil y Comercial}
\newcommand{\titulo}{El Proceso de Desalojo: Tutela Anticipada, Cautela y Condena}
\newcommand{\autor}{Isaac Edgar Camacho Ocampo}
\newcommand{\anio}{2026}

% =========================================================
% PAQUETES
% =========================================================
\usepackage[utf8]{inputenc}
\usepackage[T1]{fontenc}
\usepackage[provide=*,spanish]{babel}

\usepackage[
    a4paper,
    top=2.5cm,
    bottom=2.5cm,
    left=3cm,
    right=2.5cm,
    headheight=15pt
]{geometry}

\usepackage{amsmath}
\usepackage{amssymb}
\usepackage{mathtools}

\usepackage{graphicx}
\usepackage{float}
\usepackage{caption}
\usepackage{subcaption}

\usepackage{booktabs}
\usepackage{array}
\usepackage{longtable}
\usepackage{tabularx}

\usepackage{enumitem}

\usepackage{xcolor}
\usepackage[most]{tcolorbox}

\usepackage{fancyhdr}
\usepackage{ragged2e}

\graphicspath{
    {./img/}
    {./graficos/}
    {./figuras/}
}

% =========================================================
% CONFIGURACIÓN GENERAL
% =========================================================
\setcounter{tocdepth}{2}
\setcounter{secnumdepth}{0}

\setlength{\parindent}{0pt}
\setlength{\parskip}{0.5em}

\setlist[itemize]{topsep=0.3em,itemsep=0.2em}
\setlist[enumerate]{topsep=0.3em,itemsep=0.2em}

\clubpenalty=10000
\widowpenalty=10000

\addto\captionsspanish{\renewcommand{\chaptername}{Bloque}}

% =========================================================
% COMANDOS MATEMÁTICOS
% =========================================================
\newcommand{\abs}[1]{\left|#1\right|}
\newcommand{\norm}[1]{\left\lVert#1\right\rVert}
\newcommand{\vect}[1]{\mathbf{#1}}
\newcommand{\deriv}[2]{\frac{d#1}{d#2}}
\newcommand{\pderiv}[2]{\frac{\partial#1}{\partial#2}}

% =========================================================
% ENTORNOS / CAJAS ACADÉMICAS
% =========================================================
\newtcolorbox{definition}[1]{
    enhanced, breakable,
    title={Definición: #1},
    fonttitle=\bfseries,
    colback=blue!3, colframe=blue!50!black,
    boxrule=0.8pt, arc=2mm
}

\newtcolorbox{important}{
    enhanced, breakable,
    title={Importante},
    fonttitle=\bfseries,
    colback=yellow!8, colframe=orange!70!black,
    boxrule=0.8pt, arc=2mm
}

\newtcolorbox{example}[1]{
    enhanced, breakable,
    title={Caso / Ejemplo: #1},
    fonttitle=\bfseries,
    colback=green!4, colframe=green!40!black,
    boxrule=0.8pt, arc=2mm
}

\newtcolorbox{formula}[1]{
    enhanced, breakable,
    title={Norma: #1},
    fonttitle=\bfseries,
    colback=purple!4, colframe=purple!50!black,
    boxrule=0.8pt, arc=2mm
}

\newtcolorbox{exercise}[1]{
    enhanced, breakable,
    title={Ejercicio: #1},
    fonttitle=\bfseries,
    colback=gray!5, colframe=gray!60!black,
    boxrule=0.8pt, arc=2mm
}

\newtcolorbox{note}{
    enhanced, breakable,
    title={Nota},
    fonttitle=\bfseries,
    colback=white, colframe=gray!50,
    boxrule=0.6pt, arc=2mm
}

% =========================================================
% ENCABEZADOS Y PIES
% =========================================================
\pagestyle{fancy}
\fancyhf{}
\fancyhead[L]{\nouppercase{\leftmark}}
\fancyhead[R]{\materia}
\fancyfoot[C]{\thepage}
\renewcommand{\headrulewidth}{0.4pt}
\renewcommand{\footrulewidth}{0pt}

\fancypagestyle{plain}{
    \fancyhf{}
    \fancyfoot[C]{\thepage}
    \renewcommand{\headrulewidth}{0pt}
}

% =========================================================
% HIPERVÍNCULOS Y METADATOS PDF
% =========================================================
\usepackage[
    colorlinks=true,
    linkcolor=blue!50!black,
    citecolor=green!40!black,
    urlcolor=blue!60!black,
    pdftitle={\titulo},
    pdfauthor={\autor},
    pdfsubject={Apuntes universitarios},
    bookmarks=true
]{hyperref}

% =========================================================
% DOCUMENTO
% =========================================================
\begin{document}

% ---------------------------------------------------------
% PORTADA
% ---------------------------------------------------------
\begin{titlepage}
\thispagestyle{empty}

\begin{center}

\vspace*{1cm}

{\large \textsc{\universidad}}

\vspace{0.2cm}

{\large \textsc{\facultad}}

\vspace{3cm}

{\Huge \textbf{\materia}}

\vspace{1cm}

{\LARGE \titulo}

\vspace{0.8cm}

\textit{Comprender el proceso es el primer paso para ejercer el derecho con criterio.}

\vspace{1.2cm}

\rule{0.70\textwidth}{0.5pt}

\vspace{0.5cm}

{\large Apuntes de estudio}

\vspace{0.5cm}

\rule{0.70\textwidth}{0.5pt}

\vfill

{\large \autor}

\vspace{0.3cm}

{\large \anio}

\end{center}

\vfill

\begin{flushleft}
\textit{Este es un modesto aporte a los estudiantes de ``Derecho Procesal Civil y Comercial'' no pretende sustituir el material oficial de la materia.}
\end{flushleft}

\end{titlepage}

% ---------------------------------------------------------
% ÍNDICE
% ---------------------------------------------------------
\tableofcontents
\clearpage

% ---------------------------------------------------------
% INTRODUCCIÓN
% ---------------------------------------------------------
\chapter*{Introducción}
\addcontentsline{toc}{chapter}{Introducción}
\markboth{Introducción}{Introducción}

Este documento reconstruye una clase práctica sobre el proceso de desalojo en el derecho procesal argentino, con eje en el artículo 684 bis del Código Procesal Civil y Comercial de la Nación: la desocupación inmediata del inmueble.

Se desarrollan, con casos reales de estudio profesional, los siguientes ejes temáticos: la diferencia entre una medida cautelar y una sentencia de condena; el funcionamiento del artículo 684 bis y sus causales (falta de pago, vencimiento de contrato e intrusión); las formas de constituir la caución (autoembargo, caución real, depósito judicial); la medida cautelar innovativa del artículo 230 y el caso ``Camacho Acosta''; la perspectiva de vulnerabilidad y de género aplicada al desalojo; los medios de impugnación (efecto devolutivo y efecto suspensivo); y el acceso a la justicia a través de un lenguaje claro.

\textbf{Utilidad profesional y humana de esta unidad.} El desalojo es uno de los procesos más frecuentes de la práctica civil: cualquier profesional, ejerza para locadores o para locatarios, va a enfrentarlo. Comprender la diferencia entre cautela y condena, saber cómo se constituye una caución y conocer las causales del artículo 684 bis evita errores que cuestan tiempo, dinero y hasta responsabilidad profesional, tal como muestra el caso del cambio de cuenta bancaria desarrollado en esta clase. Además, el eje de la clase no es solamente técnico: enseña a escuchar al cliente antes de actuar, a considerar la vulnerabilidad de ambas partes del conflicto ---no solo de quien ocupa el inmueble, sino también de quien vive del alquiler--- y a preferir la solución del conflicto por sobre la resolución impuesta.

\textit{Hilo conductor de la clase: el proceso de desalojo puede pensarse como el trámite de un vecino que deja de pagar el alquiler. Primero hay que distinguir si el juez está tomando una medida urgente o dictando la sentencia final (Bloque I). Para tomar esa medida urgente, el juez exige una garantía ---como pedir una fianza antes de entregar la llave--- que se resuelve con herramientas concretas de la práctica (Bloque II). No toda persona que ocupa un inmueble es igual: hay causales distintas y personas más vulnerables que otras, y esas decisiones pueden apelarse (Bloque III). Y detrás de todo el trámite hay un objetivo mayor: que la justicia llegue a tiempo y se entienda, incluso para quien nunca pisó una facultad de derecho (Bloque IV).}

La clase se introduce, antes de entrar en el contenido técnico, con una referencia al arte como puerta de entrada al proceso judicial.

\begin{example}{la analogía de ``Las dos Fridas''}
El cuadro de Frida Kahlo ``Las dos Fridas'' muestra a la misma persona representada de dos formas distintas: un autorretrato doble. La clase lo utiliza como metáfora del expediente judicial: la demanda y la contestación de demanda narran los mismos hechos desde dos perspectivas diferentes, sin que eso implique necesariamente que una de las partes mienta. En el proceso judicial, la parte actora hace afirmaciones relativas a los hechos de la conflictiva, y la parte demandada puede hacer otra afirmación sobre esos mismos hechos: es la misma realidad, representada de manera distinta.
\end{example}

% ===========================================================
% BLOQUE I
% ===========================================================
\chapter{Cautela y Condena}
\markboth{Bloque I. Cautela y Condena}{Bloque I. Cautela y Condena}

\section{Unidad 1. Diferencia entre cautela y condena}

La resolución de condena es una resolución definitiva, que luego podrá tener o no la vía de la etapa recursiva, pero que resulta definitiva. La cautela, en cambio, es una resolución provisoria, que puede llegar a quedar firme, modificarse o revocarse.

Esta distinción resulta relevante porque, en ciertos casos, el juez o la jueza, en el plano cautelar, resuelve lo mismo que resolvería al finalizar el proceso mediante la condena. En el proceso de desalojo, el juez, con una cautela, puede resolver la desocupación inmediata del inmueble; y la resolución final del proceso ---la condena--- ordena, precisamente, la desocupación del inmueble. Se trata de la misma consecuencia práctica, situada en dos planos distintos: uno provisorio (cautelar) y otro definitivo (la condena). Que el juez adelante una decisión cautelar no debe interpretarse como prejuzgamiento.

\begin{important}
Cautela y condena no son lo mismo, aunque en el proceso de desalojo puedan ordenar la misma consecuencia práctica (la desocupación del inmueble).

La cautela es provisoria: puede confirmarse, modificarse o revocarse según lo que se pruebe en el proceso. La condena es definitiva.

Que el juez adelante una decisión cautelar no implica prejuzgamiento, porque puede modificar su criterio al dictar sentencia definitiva.
\end{important}

\section{Unidad 2. El sistema del artículo 684 bis}

El artículo 684 bis del Código Procesal fue incorporado al Código Procesal de 1967 para los procesos de desalojo. Su función consiste en adelantarse en la línea del proceso: adelantarse a la condena y llegar antes con un lanzamiento.

\begin{formula}{artículo 684 bis, Código Procesal Civil y Comercial de la Nación}
Incorporado al Código Procesal ---sancionado en 1967--- para los procesos de desalojo, habilita al juez a disponer la desocupación inmediata del inmueble en forma anticipada a la sentencia definitiva, una vez vencido el plazo para contestar la demanda y cumplidas las cargas de citación y emplazamiento, previa fijación de una caución.
\end{formula}

\section{Unidad 3. Etapas del proceso y superposición cautelar}

El proceso de desalojo se despliega, como todo proceso, a lo largo de una serie de etapas: la etapa introductoria (con el proceso ordinario del artículo 360 como referencia), la etapa probatoria, la etapa decisoria, la etapa recursiva y la etapa de ejecución, sin perjuicio del lugar que pueda ocupar la mediación previa.

En la etapa introductoria se producen los actos procesales de postulación: demanda y contestación de demanda. La parte actora ---el locador--- tiene como objeto del proceso, en la demanda, el lanzamiento: el objetivo final del proceso es recuperar el uso del inmueble. Sin embargo, mediante el artículo 684 bis, puede solicitar una desocupación inmediata anticipándose a las etapas del proceso.

Por el principio de defensa en juicio, siempre existe traslado a la otra parte. Vencido el plazo para contestar la demanda y cumplidas las cargas de citación y emplazamiento, el juez o la jueza decidirá si exige caución para dar lugar al 684 bis: esto significa llevar adelante la cautela del lanzamiento en una etapa anterior a la que naturalmente le correspondería, adelantando etapas hasta llegar a la resolución de condena.

En términos generales se exige caución. La razón es la siguiente: si la parte actora obtuvo el uso y goce del inmueble en forma anticipada, el proceso debe continuar; y si, al dictar el fallo definitivo, la jueza o el juez determina que se pidió el lanzamiento de una persona que pagaba, que estaba al día y que no cambió el destino del inmueble, quien solicitó la medida deberá responder, con la caución exigida, por todos los daños y perjuicios ocasionados. Esa es la razón de ser de la caución.

El 684 bis habilita la desocupación por falta de pago y por vencimiento del contrato, y también ---como se desarrolla en el Bloque III--- por intrusión.

% ===========================================================
% BLOQUE II
% ===========================================================
\chapter{La Caución y la Práctica Profesional}
\markboth{Bloque II. La Caución y la Práctica Profesional}{Bloque II. La Caución y la Práctica Profesional}

\section{Unidad 4. Modalidades de caución}

La forma de fijar las cauciones no está reglada por la ley, lo que genera criterios diversos según cada juzgado.

\begin{example}{modalidades de caución en la práctica profesional}
\textbf{Autoembargo:} el propio locador ofrece su inmueble en garantía. Cuando termina el proceso, se levanta el embargo; es la forma de constituir la caución sin necesidad de disponer de dinero en efectivo.

\textbf{Caución real sobre otro bien:} si el locador no desea comprometer el inmueble en disputa, puede ofrecer otro bien de su propiedad (por ejemplo, un terreno).

\textbf{Depósito dinerario judicial:} consiste en depositar una suma de dinero en el banco de los tribunales. El monto varía según el criterio de cada juzgado: algunos piden el equivalente a seis meses de alquiler, otros a tres meses, y otros hasta casi dos años, lo que suele generar resistencia por parte del cliente.
\end{example}

Cuando el inmueble en disputa es el que se ofrece en autoembargo, ante cualquier resolución desfavorable ---es decir, si no se tuvo razón en la sentencia--- se responderá por los daños con ese inmueble. El daño puede pagarse en dinero; si no se abona, el inmueble se ejecuta.

\begin{important}
La caución no tiene un monto ni una forma reglada por la ley: depende del criterio de cada juzgado.

Las alternativas más usadas en la práctica son el autoembargo del propio inmueble en disputa, la caución real sobre otro bien y el depósito dinerario judicial.
\end{important}

\section{Unidad 5. Caso práctico: el cambio de banco}

\begin{example}{el cambio de cuenta bancaria y la falta de pago}
Un locador de edad avanzada, por desconfianza hacia su banco de siempre, abrió una cuenta en otra entidad y le avisó a su inquilina, de manera informal, que a partir de ese momento debía depositar los alquileres en la nueva cuenta. No se firmó ningún documento que formalizara el cambio de domicilio de pago.

La inquilina, de buena fe, comenzó a depositar en la nueva cuenta y le enviaba capturas de los comprobantes por WhatsApp al locador, quien ---según su relato posterior--- no las revisaba.

El locador inició un proceso de desalojo por falta de pago, sin saber que la inquilina efectivamente pagaba, aunque en la cuenta nueva. Su abogado, sin ese dato, no podía adivinar dónde pagaba la inquilina.

La inquilina, ya asesorada, aportó el extracto bancario y las capturas de WhatsApp. Se contactó directamente al abogado de la contraparte para proponer una reunión y esclarecer la situación antes de continuar el litigio.

\textbf{Resultado:} se acreditó que la inquilina había pagado siempre, solo que a una cuenta distinta por una comunicación informal y no escrita. El proceso, iniciado sin causa real, se desactivó sin necesidad de ejecutar la caución.
\end{example}

La enseñanza del caso es metodológica: hay que hacer hablar al cliente y escucharlo, sin dirigir la entrevista, dejándolo contar su versión de los hechos.

\begin{formula}{recomendación de redacción contractual}
Ante cualquier cambio del domicilio de pago o de la cuenta bancaria designada en el contrato de locación, la comunicación debe formalizarse por escrito. Se recomienda incorporar en el contrato una cláusula de comunicación periódica sobre el estado del inmueble y del pago, para dejar constancia documental que evite conflictos como el descripto.
\end{formula}

\begin{important}
Escuchar al cliente es central en la práctica profesional: muchos conflictos derivados de problemas de comunicación ---y no de mala fe--- terminan en procesos judiciales que podrían evitarse.

Un cambio de domicilio de pago no formalizado por escrito puede generar un proceso de desalojo infundado, con costos y desgaste para ambas partes.
\end{important}

\section{Unidad 6. La medida cautelar innovativa (artículo 230) y el caso Camacho Acosta}

La medida innovativa del artículo 230 constituye un antecedente conceptual necesario para comprender el 684 bis.

\begin{example}{leading case ``Camacho Acosta'' (Fallos 320:1633)}
Un trabajador pierde uno de sus miembros superiores a raíz de un accidente de trabajo. Mientras avanza el proceso judicial, el muñón se va endureciendo, por lo que necesita colocarse una prótesis, cara y de difícil consecución, con urgencia, sin poder esperar el dictado de la sentencia definitiva de condena.

La Corte ordenó, por vía de una medida cautelar innovativa, que el empleador entregara la prótesis antes de que finalizara el proceso. El caso marcó una diferencia relevante: los juzgados deben, ante situaciones de esta naturaleza, proteger primero y conocer después, protegiendo primero a la parte que atraviesa la urgencia y analizando luego si existía o no derecho.
\end{example}

\begin{formula}{artículo 230, Código Procesal Civil y Comercial de la Nación --- medida cautelar innovativa}
Procede cuando se acreditan los presupuestos generales de toda medida cautelar: verosimilitud del derecho (la apariencia del derecho invocado) y peligro en la demora (que esperar el trámite ordinario del proceso torne ineficaz la protección). Requiere, además, contracautela o caución. La doctrina de Calamandrei la identifica como el sistema matriz de todas las medidas cautelares: ante cualquier situación de peligro que no encuadre en una figura cautelar específica (embargo, inhibición general de bienes, prohibición de contratar), puede recurrirse a esta vía general.
\end{formula}

El artículo 230 constituye el sistema matriz de las medidas cautelares. Antes de que existiera el artículo 684 bis, este sistema general ya existía, y algunos profesionales más avezados podían recurrir a él para pedir un lanzamiento anticipado en un desalojo, aunque no era una práctica cotidiana.

La medida innovativa admite dos variantes: la medida de no innovar, mediante la cual se solicita al juez que mantenga la situación tal como está, para que no se modifique en perjuicio de la parte que la solicita; y la medida innovativa propiamente dicha, mediante la cual se solicita al juez que modifique la situación actual, porque esa situación está causando un daño. Esta distinción, explicada aquí aplicada al desalojo, resulta aplicable a todo proceso.

\begin{important}
El artículo 230 es el sistema general (``matriz'') de toda medida cautelar; el artículo 684 bis es un subsistema especial que el legislador incorporó específicamente para el proceso de desalojo en 1967.

``Camacho Acosta'' resume la lógica de toda tutela anticipada: primero proteger, después analizar el derecho de fondo.

El 684 bis, a diferencia de una cautelar del artículo 230, exige que la otra parte ya haya sido notificada de la demanda y que haya vencido su plazo para contestar: no se dicta ``a espaldas'' del demandado, como sí puede ocurrir con un embargo cautelar tradicional.
\end{important}

% ===========================================================
% BLOQUE III
% ===========================================================
\chapter{Causales, Vulnerabilidad y Jurisprudencia}
\markboth{Bloque III. Causales, Vulnerabilidad y Jurisprudencia}{Bloque III. Causales, Vulnerabilidad y Jurisprudencia}

\section{Unidad 7. Causales de procedencia del 684 bis}

\begin{formula}{causales del artículo 684 bis}
1) Falta de pago del canon locativo. 2) Vencimiento del plazo contractual. 3) Intrusión: ocupación del inmueble sin consentimiento, bajo ninguna forma, de quien tiene derecho a otorgar su uso y goce.
\end{formula}

Respecto del concepto de ``intruso'', corresponde precisar que no toda persona que ocupa un inmueble sin contrato vigente es un intruso: si existe un contrato del año 2000, se trata de una locación continuada a la que nunca se le renovó el contrato, y no de una intrusión. El intruso es quien ocupa un inmueble sin el consentimiento, bajo ninguna forma, de la parte que otorga el uso y goce: quien lo ocupa sin derecho alguno.

Cuando se invoca el abandono del inmueble, ese abandono debe comprobarse. Para ello corresponde solicitar formalmente al juez que libre un mandamiento de constatación, con oficial de justicia de la zona ---y, eventualmente, cerrajero, si no queda otra forma de acceder---, para corroborar que en el interior no hay nada, o que hay bienes en mal estado. No corresponde autorizar al cliente a ingresar por su cuenta, por más que sostenga que ya lo hizo: debe tenerse especial cuidado con el abandono, porque el contrato puede seguir vigente y la persona puede regresar.

\begin{important}
El 684 bis solo procede por tres causales taxativas: falta de pago, vencimiento del contrato e intrusión.

``Intruso'' tiene un significado técnico preciso: ocupación sin consentimiento alguno del titular del derecho de uso y goce. Un contrato vencido o una locación continuada no convierten al ocupante en intruso.
\end{important}

\section{Unidad 8. Perspectiva de vulnerabilidad y de género}

La vulnerabilidad puede presentarse de ambos lados de la relación locativa: el cliente locador también puede ser la parte vulnerable, no solo el locatario. Existen numerosos casos de personas sin familiares, o con familiares sin recursos, que viven del alquiler de un departamento o un local pequeño, junto con una jubilación mínima. No todos los locadores son vulnerables, ni todos los locatarios lo son: existen locatarios que pagan cánones locativos muy altos, incluso en dólares, por lo que deben analizarse las circunstancias de cada caso.

\begin{example}{el lanzamiento demorado por la presencia de menores}
En un proceso con sentencia firme de desalojo, ya en la etapa de lanzamiento, se presentó una familia con menores a cargo alegando ocupar el inmueble, sin poder acreditarlo con facilidad.

El juzgado exigió la partida de nacimiento que probara el vínculo con los menores y, una vez acreditado, ofició a organismos estatales ---nacionales y municipales--- para determinar qué ayuda o reubicación podía ofrecerse a ese grupo familiar, condicionando el lanzamiento a esa respuesta.

La respuesta de los organismos oficiados fue, en la práctica, un desentendimiento generalizado: cada área remitía la competencia a otra, sin que ninguna asumiera una solución concreta.

En el mismo expediente se detectaron indicios de que el ocupante original había cedido el inmueble a esta familia deliberadamente para dilatar el lanzamiento, lo que fue calificado como un posible abuso del proceso, pasible de ser sancionado por el juez.
\end{example}

A nivel de la Ciudad de Buenos Aires existe un mecanismo mediante el cual, a través de los Centros de Gestión y Participación (CGP) de las comunas, la persona interesada lleva su contrato de locación, acredita que tiene trabajo, y el Gobierno de la Ciudad puede otorgar una subvención para evitar que la familia quede en la calle, o para buscar una reubicación. Este mecanismo requiere estar radicado en la Ciudad de Buenos Aires; en otros distritos, como La Matanza, esa herramienta no está tan desarrollada.

\begin{important}
En el tema del desalojo suele mirarse el lado interno de la puerta y no el lado externo: la vulnerabilidad debe analizarse en ambas partes del conflicto, no solo en quien ocupa el inmueble.

Los juzgados pueden condicionar la ejecución de un lanzamiento a la respuesta de organismos estatales cuando hay menores u otras personas en situación de vulnerabilidad involucradas.

El uso de la ocupación con fines dilatorios ---para retrasar artificialmente un lanzamiento con sentencia firme--- puede constituir un abuso del proceso sancionable.
\end{important}

Existen además precedentes de perspectiva de género y de familia en materia de desalojo, entre ellos un fallo de la provincia de Corrientes que rechazó un desalojo entre familiares ---dirigido contra una cuñada--- por entender que, detrás de un planteo patrimonial, existía en realidad un conflicto propio del fuero de familia y no del fuero civil patrimonial.

\begin{note}
Cuando el conflicto que motiva un pedido de desalojo entre parientes tiene, en el fondo, una raíz familiar y no patrimonial, corresponde encauzarlo por el fuero de familia y no por el fuero civil patrimonial.
\end{note}

\section{Unidad 9. Medios de impugnación y jurisprudencia}

La resolución que dispone la desocupación inmediata es impugnable. Ante la notificación de esa resolución, la parte afectada cuenta con un método impugnativo: una etapa recursiva para interponer el recurso correspondiente.

\begin{important}
\textbf{Efecto devolutivo:} la orden se cumple igual, aunque el recurso esté pendiente de resolución por la Cámara.

\textbf{Efecto suspensivo:} el juez suspende la ejecución de la resolución impugnada hasta que el superior ---la Cámara--- resuelva el recurso.
\end{important}

El recurso se interpone una vez notificada la desocupación inmediata resuelta. Conforme a la jurisprudencia citada en la clase, la Sala H le dio efecto devolutivo, es decir, la resolución impugnada se cumple igual, aunque se haya recurrido y se pretenda que un superior la revise.

\begin{example}{jurisprudencia citada sobre el artículo 684 bis}
\textbf{Sala D de la Cámara Nacional Civil:} confirmó la procedencia del 684 bis en un caso en el que ya se había concretado el lanzamiento con auxilio de la fuerza pública, resolviendo continuar adelante pese a la oposición planteada.

\textbf{Fallo del departamento judicial de La Plata:} un precedente provincial que también reconoce la procedencia de esta figura, mostrando que tanto los fallos nacionales como los provinciales acogen el instituto.
\end{example}

\begin{important}
La resolución de desocupación inmediata es recurrible, pero en la mayoría de los precedentes citados el recurso se concede con efecto devolutivo: el lanzamiento se ejecuta igual mientras la Cámara resuelve.

El efecto suspensivo es la excepción: exige al juez de primera instancia detener la ejecución hasta que el superior se expida.
\end{important}

% ===========================================================
% BLOQUE IV
% ===========================================================
\chapter{Acceso a la Justicia y Reflexión Final}
\markboth{Bloque IV. Acceso a la Justicia y Reflexión Final}{Bloque IV. Acceso a la Justicia y Reflexión Final}

\section{Unidad 10. Lenguaje claro y acceso a la justicia}

\begin{example}{la cédula que no se entendió a tiempo}
Una señora recibió la cédula de notificación de una demanda de desalojo, pero no comprendió la urgencia del documento y no lo envió por foto a su hijo, quien estaba de viaje con su nieto en un viaje de egresados de séptimo grado. No quiso interrumpir el viaje del nieto.

El domicilio estaba correctamente notificado ---el hijo, soltero, figuraba con su documento en ese domicilio---, por lo que la notificación era formalmente válida. Sin embargo, la falta de comprensión del documento por parte de una persona sin formación jurídica generó la pérdida de un plazo procesal relevante.
\end{example}

La gente sin formación jurídica, con frecuencia, no advierte que en esa notificación ---percibida como un documento hostil, redactado además en un lenguaje técnico--- hay una fecha fijada por el oficial notificador, y las expresiones que en el ámbito jurídico se consideran claras no siempre lo son para quien no conoce esa terminología.

En algunos casos, la redacción compleja de los autos judiciales resulta a propósito, y esa práctica es incorrecta, porque toda persona tiene derecho ---así lo garantizan la Constitución Nacional, la Convención Americana sobre Derechos Humanos y el resto del bloque de constitucionalidad--- a recibir información clara sobre los caminos procesales que puede tomar ante un juicio.

\begin{important}
El derecho a la información procesal, en un lenguaje claro y comprensible, es un derecho constitucional y convencional de todo ciudadano, no un tecnicismo reservado a quienes estudiaron derecho.

Una notificación formalmente válida puede, en la práctica, fracasar en su función si el destinatario no comprende su contenido ni su urgencia.
\end{important}

\section{Unidad 11. Flexibilización de las formas y tutela judicial efectiva}

Existen fallos provinciales, en general más innovadores en esta materia, en los que se ha permitido contestar demanda fuera del plazo procesal. En un caso relatado en clase, una jueza de familia dejó contestar la demanda fuera de término, dejando en claro esa circunstancia, pero aclarando que ello no significaba adelantar su resolución: solo se buscaba escuchar por qué no se había contestado dentro de los cinco días del proceso sumarísimo, qué prueba se ofrecía, y luego analizar si correspondía proveerla o no en el marco de la jurisdicción.

Esta flexibilización, orientada a llegar a la verdad jurídica material, entra en tensión con las reglas del sistema de preclusión y puede abrir la puerta a la inseguridad jurídica. Sin perjuicio de las diferencias entre provincias y de las circunstancias de tiempo y lugar ---no es lo mismo litigar en una gran ciudad que en un pueblo alejado---, el fuero de familia admite, en general, flexibilidades que no existen en el fuero patrimonial.

\begin{important}
La flexibilización de los plazos procesales es una tensión real entre la búsqueda de la verdad jurídica material y la seguridad jurídica que da el sistema de preclusión.

El fuero de familia admite, en la práctica de algunas provincias, mayores flexibilidades formales que el fuero civil patrimonial.
\end{important}

\section{Unidad 12. Derecho y arte: síntesis de la clase}

La clase se cierra con la imagen de la balanza de la justicia y un paralelo constante: la búsqueda de soluciones de conflicto, en lugar de resoluciones impuestas por la jueza o el juez, a lo que se suma la necesidad de agilizar los procesos en aras de la tutela judicial efectiva. Esta agilización resulta fundamental, por ejemplo, para el jubilado que vive con un ingreso mínimo y necesita recuperar rápidamente su pequeño inmueble para poder volver a alquilarlo.

La referencia a Frida Kahlo introducida al comienzo de la clase no fue meramente anecdótica: en el mundo, y particularmente en Europa, se observa una intersección creciente entre derecho y arte, como una forma de salir de la saturación de los enfoques puramente técnicos, relacionando obras y referencias culturales con el contenido de la materia.

\begin{important}
El proceso de desalojo, más allá de su tecnicismo, es un espacio donde conviene priorizar la solución del conflicto por sobre la resolución judicial impuesta.

La tutela judicial efectiva exige agilizar los procesos, especialmente cuando están en juego personas en situación de vulnerabilidad económica, como jubilados que dependen del alquiler de un pequeño inmueble.

El vínculo entre derecho y arte ---ejemplificado con ``Las dos Fridas''--- funciona como una herramienta pedagógica para pensar que una misma realidad admite más de una representación legítima dentro del proceso.
\end{important}

% ===========================================================
% RESUMEN CORNELL
% ===========================================================
\clearpage
\chapter*{Resumen Cornell}
\addcontentsline{toc}{chapter}{Resumen Cornell}
\markboth{Resumen Cornell}{Resumen Cornell}

Cuadro de repaso con preguntas disparadoras y su desarrollo, más una síntesis integradora final de toda la clase.

\begin{longtable}{
    >{\RaggedRight\arraybackslash}p{0.30\textwidth}
    >{\RaggedRight\arraybackslash}p{0.60\textwidth}
}
\toprule
\textbf{Preguntas / palabras clave} & \textbf{Notas / respuestas} \\
\midrule
\endfirsthead
\toprule
\textbf{Preguntas / palabras clave} & \textbf{Notas / respuestas} \\
\midrule
\endhead

\textbf{¿Cautela vs. condena?} &
La cautela es una resolución provisoria (puede confirmarse, modificarse o revocarse); la condena es la resolución definitiva del proceso. En el desalojo, ambas pueden ordenar lo mismo ---la desocupación del inmueble--- sin que eso implique prejuzgamiento. \\
\addlinespace
\cmidrule{1-2}

\textbf{Artículo 684 bis CPCCN} &
Incorporado en 1967, permite al juez ordenar la desocupación inmediata del inmueble antes de la sentencia definitiva, una vez vencido el plazo de contestación de demanda y cumplidas las cargas de notificación, previa caución. \\
\addlinespace
\cmidrule{1-2}

\textbf{¿Por qué se exige caución?} &
Porque si quien pidió la desocupación anticipada no tenía razón, debe responder con esa garantía por los daños y perjuicios ocasionados a la otra parte. \\
\addlinespace
\cmidrule{1-2}

\textbf{Formas de constituir la caución} &
Autoembargo del propio inmueble en disputa, caución real sobre otro bien del locador, o depósito dinerario en el banco de tribunales (el monto varía según cada juzgado). \\
\addlinespace
\cmidrule{1-2}

\textbf{Caso del cambio de banco: lección} &
Un locador cambió de cuenta sin formalizarlo por escrito; la inquilina pagaba igual, a la cuenta nueva. El desalojo por falta de pago se inició sin causa real. Lección: escuchar al cliente y dejar todo cambio de domicilio de pago por escrito. \\
\addlinespace
\cmidrule{1-2}

\textbf{Artículo 230 CPCCN y caso Camacho Acosta} &
El 230 es el sistema matriz de toda medida cautelar (verosimilitud del derecho + peligro en la demora + contracautela). Camacho Acosta es el leading case que impone ``proteger primero, conocer después'' ante una urgencia extrema. \\
\addlinespace
\cmidrule{1-2}

\textbf{Tres causales del 684 bis} &
Falta de pago, vencimiento del contrato e intrusión. ``Intruso'' es quien ocupa sin consentimiento alguno de la parte con derecho a otorgar el uso y goce; un contrato vencido no convierte al ocupante en intruso. \\
\addlinespace
\cmidrule{1-2}

\textbf{Perspectiva de vulnerabilidad} &
Hay que mirar los dos lados de la puerta: tanto quien ocupa el inmueble como el locador pueden ser la parte vulnerable (por ejemplo, un jubilado que vive del alquiler de un pequeño inmueble). \\
\addlinespace
\cmidrule{1-2}

\textbf{Efecto devolutivo vs. suspensivo} &
Devolutivo: la resolución se ejecuta igual mientras se resuelve el recurso (regla general en los fallos citados). Suspensivo: se detiene la ejecución hasta que el superior resuelva. \\
\addlinespace
\cmidrule{1-2}

\textbf{Lenguaje claro y acceso a la justicia} &
El derecho a comprender los caminos procesales disponibles es un derecho constitucional y convencional. Una notificación válida puede fracasar en la práctica si el destinatario no entiende su urgencia. \\
\addlinespace
\cmidrule{1-2}

\textbf{Síntesis final de la clase} &
Priorizar la solución del conflicto por sobre la resolución impuesta, agilizar los procesos en pos de la tutela judicial efectiva, y recordar ---con la imagen de ``Las dos Fridas''--- que una misma realidad admite más de una representación legítima en el proceso. \\
\addlinespace
\midrule

\multicolumn{2}{p{0.94\textwidth}}{
\textbf{Resumen:} El proceso de desalojo argentino, articulado alrededor del artículo 684 bis del Código Procesal Civil y Comercial de la Nación, ilustra con particular claridad la tensión estructural entre dos valores del proceso civil: la necesidad de una tutela judicial rápida y efectiva, y la garantía de defensa en juicio de quien resulta afectado por una medida antes de la sentencia definitiva. La distinción entre cautela (provisoria) y condena (definitiva) resulta la clave conceptual de toda la clase, y se combina con el sistema general de las medidas cautelares del artículo 230 ---cuya lógica de ``proteger primero, conocer después'' ilustra el caso Camacho Acosta--- para mostrar que el 684 bis no es una cautelar más, sino un subsistema especial en el que, a diferencia de una cautelar clásica, la parte demandada ya fue notificada de la demanda y tuvo su oportunidad de defensa antes de que se ordene la desocupación anticipada. Sobre esa base técnica, la clase suma una capa de criterio profesional y humano: la caución no está reglada y exige creatividad práctica (autoembargo, caución real, depósito judicial); la vulnerabilidad debe mirarse de los dos lados de la puerta, porque tanto quien ocupa el inmueble como quien vive de alquilarlo pueden ser la parte débil del conflicto; y el acceso efectivo a la justicia depende, en última instancia, de que las personas comunes puedan comprender el lenguaje de la notificación que reciben. El hilo que atraviesa toda la clase ---remarcado con la imagen de ``Las dos Fridas'' de Frida Kahlo--- es que el proceso judicial es, ante todo, una construcción de sentido sobre una misma realidad contada desde perspectivas distintas, y que la función de quien ejerce el derecho no es únicamente aplicar la letra de la norma, sino buscar, siempre que sea posible, una solución del conflicto antes que una resolución impuesta.
} \\
\bottomrule
\end{longtable}

\end{document}
